% Apéndice: material de clase y notebooks
\section{Material de clase y notebooks}\label{cap:21}\label{ap:notebooks}

\subsection{Material publicado en el campus (cursada 2025)}

Copia local en \file{Proyectos / Agrimensura / Sistemas de Coordenadas / Bibliografia / material de
clase 2C2025} de Dropbox.

\begin{itemize}
  \item \file{Rotaciones_No_Conmutan.xlsx} (Planilla Clase 2025.2.2, clase 2): planilla
  para ``experimentar con la conmutatividad de las matrices de rotación y con la
  construcción de una rotación según eje-ángulo (por repetición de rotaciones pequeñas y por
  exponenciación)''. \texttt{MMULT} es la función de Excel para multiplicar matrices.
  \item \file{Ejercicio Tl_Afin_Homogenea.xlsx} (Planilla para experimentar
  transformaciones, clase 3): afinidades en \(\R^2\) y en homogéneas; composición y orden.
  \item \file{Algun_ejemplo_de_rotacion.wxmx} (Maxima 2025.2.2, clase 2): notebook de
  wxMaxima que ``experimenta nuevamente el tema de la conmutatividad y aparece una yapa
  relativa a autovectores''.
  \item \file{Formas_Lineales.wxmx} (clase 4): el ejemplo \(\varphi(A\vec v)=(\varphi A)(\vec v)\)
  trabajado en Maxima; es el de la sección~\ref{sec:formas-lineales}.
  \item \file{Ejemplo_para_pensar.py} (+ \texttt{.jpeg}; Ejemplo Clase 2025.2.2): ejemplo
  en Python de transformación afín ``que vamos a trabajar en papel en clase [\ldots]
  piensen qué hace''.
\end{itemize}

\subsection{Unidades interactivas (lecciones marimo)}

Las 19 unidades (00 a 14) son el hilo conductor del sitio de la materia: lecciones marimo +
Python escritas por mí como alumno, construidas ladrillo a ladrillo y publicadas como
páginas WASM (corren en el navegador, sin instalar nada):

\begin{center}\url{https://gusvarela-fiuba.github.io/Sistemas-de-Coordenadas/}\end{center}

El sitio tiene tres capas: ese hilo conductor (la portada), el mapa de conceptos
(herramienta de repaso) y este apunte formal en \LaTeX{}. El fuente está en
\file{unidades/} del repositorio; se abren localmente con \texttt{marimo edit}
(entorno \texttt{uv}). Cada unidad enlaza a la sección de este apunte que desarrolla su
tema (destino \texttt{cap:NN} del PDF), y cada una cierra con ejercicios con hueco y
verificación automática. La columna ``clase'' es la del esqueleto de 14 clases de la
cursada 2025 (sección~\ref{sec:hilo}). La Unidad 0 va antes de la clase 1: es el
encuadre de la subsección~\ref{sec:desaprender} y de la brújula de la
sección~\ref{sec:puntos-formas-grillas}, no una clase. Las unidades 11 a 14 son
\emph{anticipadas}: la clase no se dictó en 2C2026 y el contenido es lectura mía sobre el
esqueleto 2025 y el apunte del elipsoide del profesor, pendiente de validar por la
cátedra.

{\small
\begin{longtable}{>{\ttfamily\raggedright\arraybackslash}p{62mm}>{\raggedright\arraybackslash}p{70mm}>{\raggedright\arraybackslash}p{16mm}}
\toprule
\textrm{\textbf{Unidad (id)}} & \textbf{Título} & \textbf{Clase}\\
\midrule
\endfirsthead
\toprule
\textrm{\textbf{Unidad (id)}} & \textbf{Título} & \textbf{Clase}\\
\midrule
\endhead
\bottomrule
\endlastfoot
00\_desaprender\_el\_cbc & Fundamentos: cuerpos, grupos, puntos, formas lineales y métrica & ---\\
01a\_punto\_quieto & Rotaciones I: el punto no se movió & 01\\
01b\_rotacion\_plano & Rotaciones II: la rotación en el plano & 01\\
01c\_rotaciones\_3d & Rotaciones III: tres ejes y el orden importa & 01, 02\\
02\_producto\_interno & Producto interno: proyección, longitud, ángulo & 01\\
03\_reflexiones\_similaridades & Reflexiones y similaridades & 02, 03\\
04\_afines\_homogeneas & Afinidades y coordenadas homogéneas & 03\\
05\_proyectivas & Proyectivas & 03p\\
06\_formas\_lineales & Formas lineales y dualidad & 04\\
07a\_cuna\_determinante & Cuña y determinante: el detector de colinealidad & 05\\
07b\_nucleo\_autovectores & Núcleo y autovectores: las direcciones que no giran & 05, 06\\
07c\_autovalores\_3d\_\allowbreak puntos\_fijos & Autovalores en 3D y puntos fijos proyectivos & 06\\
08\_cambio\_de\_base & Cambio de base y diagonalización & 08\\
09\_cambio\_de\_base\_svd & Cambio de base, SVD e inversas (con el DEM real de la
sección~\ref{sec:svd}) & 07, 10\\
10\_jacobiano\_primera\_forma & Jacobiano y primera forma fundamental & 09\\
11\_cuatro\_problemas & Cuatro problemas: regresión, círculos, Procrustes, regresión total
(anticipada) & 11\\
12\_curvatura & Curvatura como forma cuadrática (anticipada) & 12\\
13\_elipsoide & Geometría del elipsoide (anticipada) & 13\\
14\_trigonometria\_esferica & Trigonometría esférica (anticipada) & 14\\
\end{longtable}
}

\paragraph{Módulo común \file{cb104.py}.} Vive en \file{unidades/}, junto a los
notebooks, y todas las unidades lo importan (\texttt{import cb104 as lib}). Sólo depende
de NumPy (y de Plotly adentro de las funciones que dibujan). Reúne la paleta del apunte
(teal, rust, gold), las rotaciones \texttt{rot2}, \texttt{rot\_x}, \texttt{rot\_y},
\texttt{rot\_z} y \texttt{rodrigues}, coordenadas homogéneas (\texttt{hom}, \texttt{dehom}),
área orientada de un polígono, potencias reales \texttt{pot\_real}, cambio de base
(\texttt{coordenadas\_en\_base}, \texttt{diagonalizar}), eje y ángulo de una rotación,
helpers de \LaTeX{} para mostrar matrices coloreadas, las capas de Plotly ``antes/después''
de la grilla, el cubo unitario y los terrenos sintéticos en 3D, y la verificación de los
ejercicios con hueco (\texttt{verificar}, \texttt{exacto}). Para el export WASM el módulo
no existe en el navegador: \file{tools/build.py} reemplaza la celda marcada
\texttt{\# cb104:inline} por una que embebe el archivo. Se prueba con
\texttt{uv run python unidades/cb104.py}.

\subsection{Archivos del profesor recibidos por el chat de la cursada 2C2026}

Copias fieles en \file{fuentes/campus/profesor/} del repositorio (recibidos el 17 de
octubre de 2026 por el grupo de Google Chat, junto con los dos PDF que este apunte
reproduce al final).

\begin{itemize}
  \item \file{2026-09-16\_pizarron\_pentagono.jpg}: foto del pizarrón del 16 de septiembre
  con la consigna del pentágono: dos dibujos del pentágono \(ABCDE\) y las cuatro preguntas
  (qué proyectivas llevan cuatro vértices a otros cuatro, cómo se calculan, dónde cae el
  vértice restante, si hay punto fijo). Es la consigna del apéndice~\ref{ap:pentagono}.
  \item \file{Proyectivas\_del\_pentagono.xlsx}: planilla del profesor (una hoja,
  \texttt{B2:U59}). Pentágono regular centrado en el origen con \(D=(0,1)\) y los demás
  vértices por rotaciones de \(72^\circ\); arma el sistema \(13\times13\) de la proyectiva
  \(3\times3\) (nueve coeficientes más un factor \(\lambda\) por punto, con el último
  coeficiente fijado en \(1\)) que deja fijos \(A\) y \(B\) e intercambia \(C\) con \(E\),
  lo resuelve con la matriz inversa y verifica dónde cae cada vértice al dividir por la
  tercera coordenada. Es la planilla que el 30 de septiembre el profesor dijo haber hecho
  la clase anterior. La lectura de la planilla es mía a partir de sus fórmulas.
  \item \file{LaSVD.wxmx}: cuaderno de wxMaxima con la SVD analítica de
  \(A=\begin{pmatrix}1&2&4\\-3&-2&3\end{pmatrix}\): \(AA^T\) y \(A^TA\), polinomio
  característico, autovectores, ortogonalidad de a pares, \(U\) y \(V^T\) normalizados,
  \(\Sigma\) con las raíces de los autovalores y la verificación \(U\Sigma V^T=A\). Es ``la
  SVD con Maxima, analíticamente'' que el profesor mencionó el 14 de octubre; rehecha en
  la sección~\ref{sec:svd}. Se abre con wxMaxima (es un zip con \texttt{content.xml}).
\end{itemize}

\subsection{Índice general de problemas y guías resueltas}

Dónde está resuelto cada tipo de problema: en este apunte, en las entregas al campus
(\file{entregas/}, alumno Gustavo Varela, padrón 56798) y en los ejercicios con hueco de
las unidades.

{\small
\begin{longtable}{>{\raggedright\arraybackslash}p{34mm}>{\raggedright\arraybackslash}p{46mm}>{\raggedright\arraybackslash}p{48mm}>{\raggedright\arraybackslash}p{22mm}}
\toprule
\textbf{Tema} & \textbf{En este apunte} & \textbf{En \file{entregas/}} & \textbf{Unidades}\\
\midrule
\endfirsthead
\toprule
\textbf{Tema} & \textbf{En este apunte} & \textbf{En \file{entregas/}} & \textbf{Unidades}\\
\midrule
\endhead
\bottomrule
\endlastfoot
Desaprender el CBC: los cuatro cambios de chip &
\S\ref{sec:desaprender}, \S\ref{sec:puntos-formas-grillas} & --- & 00\\
Rotaciones 2D y 3D, eje y ángulo, exponencial &
\S\ref{cap:07}, ap.~\ref{ap:rot3d}, ap.~\ref{ap:infinitesimales} & --- & 01a, 01b, 01c\\
Guía de reflexiones, traslaciones, rotaciones, homotecias y similaridades (teoremas 1 a 6) &
enunciados en \S\ref{cap:08}, guía resuelta completa en \S\ref{cap:09} & --- & 02, 03\\
Afinidades: los seis ejercicios de la clase del 9/9 &
resolución en ap.~\ref{cap:17}; composición en ap.~\ref{ap:composicion-afin} &
\file{Afinidades\_Ejercicios\_Varela\_56798.pdf} (fuente y verificación NumPy en
\file{afinidades/}); extra no pedido: \file{Afinidades\_Triangulo\_Varela\_56798.pdf},
llevar un triángulo sobre otro con afinidades axiales & 04\\
Proyectivas del pentágono (consigna del 16/9, aclaración del 18/9) &
cuentas y script en ap.~\ref{ap:pentagono}; teoría en \S\ref{cap:12} &
\file{Pentagono\_Proyectivas\_Varela\_56798.pdf} (20 páginas: pentágono regular, tres
casos y composiciones, más el mismo problema sobre un pentágono irregular con enteros y
la parte conceptual; fuente, notebook marimo y verificación en \file{pentagono/}) &
05, 07c\\
Productos, determinante, núcleo, autovectores por producto cruz &
ejemplos resueltos en \S\ref{cap:11} (reflexión, rotación compuesta) & --- & 07a, 07b, 07c\\
Formas lineales, cambio de base, diagonalización &
\S\ref{sec:formas-lineales} (forma arrastrada), \S\ref{sec:cambio-base} (ejercicio del
30/9, \(A^5\), \(A^t\)) & --- & 06, 08\\
SVD, inversas, pseudoinversa &
\S\ref{sec:svd} (dos SVD a mano, DEM), \S\ref{sec:pseudoinversa} (recta por tres puntos)
& --- & 09\\
Los cuatro problemas: regresión, círculos, Procrustes, regresión total &
guía resuelta en \S\ref{sec:cuatro-problemas} & --- & 11\\
Jacobiano, primera forma, curvatura &
\S\ref{sec:metrica}, \S\ref{sec:curvatura} (silla, esfera) & --- & 10, 12\\
Elipsoide, arco de meridiano, proyecciones &
\S\ref{sec:elipsoide} (GRS80 en \(45^\circ\) y \(-34{,}6^\circ\)),
\S\ref{sec:arco-meridiano} (serie y cinco proyecciones) & --- & 13\\
Trigonometría esférica, exceso &
\S\ref{sec:trigonometria-esferica} (octante, tres ciudades) & --- & 14\\
\end{longtable}
}
